\section{IA Explicável (XAI) e Governança Epistemológica}
\label{sec:ia-xai}

A propensão histórica da inteligência artificial médica à opacidade inferencial — os chamados modelos de "caixa preta" (\textit{black box}) — colide frontalmente com a epistemologia clínica cartesiana e com arcabouços ético-legais rigorosos. Na práxis odontológica, em que uma classificação equivocada ou espúria pode deflagrar cirurgias iatrogênicas ou a exodontia radical de elementos hígidos, a adoção de IA Explicável (\textit{Explainable AI} --- XAI) deixou de ser um adereço cosmético de \textit{software} para se instaurar como um vetor de auditoria obrigatória \cite{ethical2026realising}.

A governança do \textit{deep learning} demanda, compulsoriamente, a decodificação da causalidade da inferência algorítmica para a linguagem da razão humana. Três metodologias matemáticas têm despontado no eixo da odontologia digital contemporânea:

\textbf{Grad-CAM (Gradient-weighted Class Activation Mapping):} Um paradigma de explicabilidade \textit{post-hoc} intrínseco à neurocomputação convolucional. Ao retropropagar o gradiente decisório das camadas latentes de volta à resolução da topologia de entrada, a técnica plasma mapas de calor (\textit{heatmaps}) termográficos evidenciando onde o algoritmo de fato alocou seu peso. Se a arquitetura classifica uma panorâmica como "alto risco para lesão endoperiodontal avançada", o Grad-CAM incendiará na imagem as regiões de trabeculado rarefato ou a perda de inserção interproximal responsável pela assertiva. A técnica exime o processo de basear sua decisão analítica em vieses espúrios ocultos, a exemplo do posicionamento equivocado do clipe radiográfico ou discrepâncias de calibração do raio-X.

\textbf{SHAP (SHapley Additive exPlanations):} Fundamentado estritamente na precisão axiomática da teoria de jogos cooperativos, o mecanismo de SHAP segmenta aditivamente os louros da classificação final por todos os vetores descritivos que alimentaram o sistema. Ao submeter os dados fenotípicos de um paciente a um modelo de avaliação de falência e rejeição peri-implantar primária, o SHAP elenca quantitativamente as influências de atributos individuais: a densidade trabecular somou $+0.4$ de propensão à falha, e o fator glicêmico basal incidiu em $+0.25$, concedendo ao odontologista a lucidez sobre em quais fatores etiológicos concentrar sua intervenção sistêmica pré-cirúrgica.

\textbf{LIME (Local Interpretable Model-agnostic Explanations):} Operando ao injetar distorções probabilísticas aos estímulos de entrada para mensurar como o algoritmo opaco tangencia suas predições nas fronteiras, o LIME forja um simulacro sub-reptício de predição estritamente linear ao redor exclusivo da predição local do indivíduo sondado. O LIME prova-se inestimável ao dissecar as razões pelas quais uma característica atípica específica deflagrou alarmes oncológicos na ausência de indicadores morfológicos óbvios, isolando o foco num escopo pontual indolor sem necessitar expor as dezenas de bilhões de parâmetros do modelo completo de Gêmeo Digital.

No âmbito legislativo do Brasil e globo afora, diplomas como a Lei Geral de Proteção de Dados (LGPD) e o Regulamento Geral sobre a Proteção de Dados (GDPR) cimentaram prerrogativas irrevogáveis sobre dados que alicerçam a autodeterminação cognitiva. A decisão de um modelo de reabsorção que define prognósticos mutilantes carece de explicação compreensível. O Gêmeo Digital Odontológico sem rastreabilidade falha epistemologicamente ao tentar impor uma sentença autocrática à equipe clínica interdisciplinar; seu papel só é legal e eticamente aceito se as inferências propostas forem, a todo instante, justificadas e abertas ao rebate pelo agente principal em saúde humana: o clínico e o pesquisador responsável, viabilizando o que plataformas pioneiras definem como Governança Comportamental e Auditoria Lógico-formal, implementada através de arquiteturas sólidas como integrações via OpenTwins \cite{infante2024integrating} e sistemas de validação de Orquestração Baseada em Regras Z3 do OpenCode.
